\section*{Network Patterns in Intrastate Conflict}

In this study, we investigate the formation and evolution of what we term the civil conflict network. This network is the product of strategic actions by both government and armed non-governmental actors. In this section, we discuss how this network emerges, what we gain by treating civil conflict as a network problem, and how civil conflict evolves and shifts when new actors enter the network.

To understand the generation of the conflict network, it is useful to start by looking at the motivations of the actors that make up this network -- the government and rebel groups. We follow \citet{collier:hoeffler:2004} in claiming that rebel actions are driven by a combination of material factors (sometimes referred to as greed) or the desire to gain additional resources, rents, or otherwise enrich the group and its leaders; and ideological factors (often called grievances) or the desire to change public policy and shift the benefits towards particular economic, social, ethnic or religious groups (see also \citet{gurr:1970, regan:norton:2005, bodea:elbadawi:2007}). Conversely, the government is motivated principally by the desire to remain in power and secondarily to enrich themselves and their winning coalition \citet{bdm:etal:2005,goemans:2000,acemoglu:robinson:2001}.
 .
These motivations provide incentives for different groups to attack other armed groups. A rebel group that is motivated by greed might attack not just the government, but other rebel groups in order to increase their abilities to control rents or lootable resources. Ideological motivation could drive groups to attack the government in hopes of either toppling the government and changing policy themselves, driving the government out of territory in order to impose policies in that territory, or simply to impose costs and obtain policy concessions. Ideologically motivated groups might also attack rebel groups that draw support from different sections of society in order to shift support away from challenger groups in a territory. The government also has a number of incentives to attack rebel groups -- attacking strong rebel groups reduces the likelihood those groups can defeat the government, but governments also attack to control resources and suppress challengers' abilities to interact violently (through predation) or nonviolently (through public goods provision) with the civilian population. 

In order to attack other groups, rebels rely on the ability to mobilize resources. They can do this by mobilizing or taxing the civilian population \citet{kalyvan:2006}, through the sale of so called ``lootable" goods like alluvial diamonds, narcotics, or tropical timber \citet{ross:2004, lujala:etal:2005, soysa:2000}, or they can gain resources from foreign supporters, generally those that are rivals of the central government \citet{gleditsch:2007, regan:2002, regan:2002a}.

Armed actors and their actions, driven by their motivations and sources of resources, constitute what we call the ``civil-conflict network." In this network the nodes are the different armed groups -- both rebel and government -- and the links are battles between armed groups. This network is not static over time. New groups are created, groups split into rivals, and groups demobilize and leave the network. Similarly, groups are not constant in their behavior over time, and groups might fight eachother in one time period, and a common enemy in another.
%CD: HIS MIGHT BE THE PLACE FOR a SHORT description of the stylized net + graphic

 % A stylized representation of a dynamic network is shown below. In Figure~\ref{fig:netPanel2}, the nodes represent armed actors and the ties represent battles between armed actors. The graphic highlights several important features relevant to our study. Each panel represents a different phase, or new year, of the conflict. In each year, new armed actors might enter or exit the conflict (i.e., an armed group could stop fighting or dissolve). In the visualization, new entrants are shaded in grey. Second, new ties can form between actors in any period, but ties can also disappear over time. For example, moving from panel 1 to panel 2 we can see that some of the linkages in the first time period no longer appear in the second time period, which indicates that two actors that fought in the first time period are not fighting in the second time period. These features demonstrate the type of information that a network conceptualization of conflict development can provide.

% \begin{figure}[ht]
% 	\centering
% 	\begin{tabular}{c|c|c}
% 		$t=1$ & $t=2$ & $t=3$ \\
% 		\includegraphics[width=.2\textwidth]{netPanel2_1.pdf} &
% 		\includegraphics[width=.2\textwidth]{netPanel2_2.pdf} &
% 		\includegraphics[width=.2\textwidth]{netPanel2_3.pdf} \\
% 	\end{tabular}
% 	\caption{Stylized network graphic depicting changing connections between actors overtime, directed linkages, and community clusters.}
% 	\label{fig:netPanel2}
% % \end{table}
% \end{figure}

\subsection{First Order Dependencies}
In the conflict network, some actors are more likely to fight than others due to their group-level characteristics. In the studies of networks, these are called first-order dependencies, and much of the study of rebel group behavior focuses on these kinds of factors.\footnote{For example \citet{weinstein:2007} on ideological groups being less violent AND OTHER STUFF.}
In particular, \emph{ceteris paribus}, certain groups are going to be more likely to \emph{attack} other groups, and certain groups are going to be more likely to \emph{be attacked.}

Many of these first order dependencies stem from a group's motivations or their resources. For example, if a group is ideologically motivated where most other groups are ideologically different, then this group is in a target rich environment, where they will have many opportunities to advance their ideology through violence. Similarly, groups that adhere to an exclusionary ideology, or an ideology that is center-seeking, as opposed to one that can be satisfied by secession, are going to be more likely to attack across the network. Another example of first order dependence is when a contextual feature of the group makes it more likely to be a target. For example, when a group possesses territory with valuable natural resources, they are more likely to be attacked by groups that are motivated by greed. 

While some of these factors can be controlled for in the context of a standard regression model, many of them are either contextual or difficult to observe (measure). This creates serious challenges for understanding who fights whom and when. If two groups fight, it might be because one of the groups is particularly aggressive, or the other group is a particularly appealing target, and failing to account for these first order dependencies will lead to inappropriately assigning credit for the confrontation, and a biased understanding of the sources of conflict. A network approach allows for us to both appropriately account for actor-level effects, or first-order dependencies, without directly measuring whether ideology or religion or resources are driving these effects. While measuring this is undoubtedly important and relevant to the study of conflict, testing hypotheses on these characteristics has already been a primary focus of existing theoretical and empirical work. Our aim is to move one step beyond these factors to explore other dimensions of the conflict system.

\subsection{Second Order Dependencies}
Networks of civil conflict are not created only by the attributes of armed groups, they are also a consequence of armed groups' actions.  CITATION have explored how armed groups affect the tendency of their targets to be violent in the future [CONDRA AND SHAPIRO/LYALL ON VICTIMIZATION, STUFF ON SPIRALS AND RETALIATION], but few studies have fully explored the consequences of armed group actions (and interactions) to explain fighting during civil wars.

One of the most important ways that conflicts evolve is through reciprocity and retaliation. If one armed group attacks another, it is more likely to be attacked itself. Some of this behavior is not causal. If a latent characteristic of groups leads actor A to attack B, it is also more likely to lead B to attack A. At the same time, dynamics of support can also encourage reciprocity. If a group does not respond to attacks against it, it risks being seen as weak and which can lead to a loss of support from civilians or foreign sponsors. For civilians, the need for retaliation can be especially stark, as one of the reasons civilians might support an armed group is for protection from other armed groups [CITE VALENTINO AND KALYVAS ON THIS].

Accounting for reciprocity in a conflict network improves our ability to understand many dynamics of civil conflict. It becomes difficult to parse the factors that cause groups to attack and to be targeted if we ignore the fact that some attacks are retaliatory in nature. More generally, understanding reciprocity during civil conflict underscores that conflicts are not static: the actions and reactions passed between groups play a crucial role in shaping the contours of civil conflict.

\subsection{Third Order Dependencies}
We posit that another form of dependency is relevant to explain conflict in the conflict network. Importantly, there are ways in which the motivations, resources, and actions of groups effect not just their targets, but also unrelated pairs of actors. In the network literature, this is termed ``third-order dependencies," where the relationship between A and B are linked with the relationship between B and C, and the relation between A and C. Two particular ways in which this instantiate are homophily and stochastic equivalence.

Armed actors' motivations often lend themselves to homophily or heterophily -- the tendency of groups that share a latent trait to interact more or less. Homophily (and heterophily) themselves cause networks to have greater (lesser) levels of transitivity and clusterability. Groups that are ideologically motivated are more likely to attack armed actors that are not ideologically similar (the same holds for ethnic or religious motivations). Similarly, groups that are motivated by greed are more likely to fight geographically proximate groups rather than distant ones. Accordingly, we should also observe that groups with similar resource bases fight eachother in an effort to monopolize access to those resources.

The second type of third order dependency that we expect to influence conflict is stochastic equivalence. Stochastic equivalence captures the tendency for actors to behave similarly in the conflict network. A very straightforward example of this is when groups are sponsored by the same foreign actor. In this case, groups are going to be unlikely to fight each other and more likely to fight the government (CITE BHAPAT AND BOND). A similar dynamic occurs when rebel groups have actual alliances (CITE SEAN'S STUFF). When groups have similar ideological motivation, they are not only unlikely to fight each other, but they are also incentivized to select similar targets. Finally, when ideological groups have secessionist aims, they could be expected to be stochastically equivalent since they would be most likely to fight the central government and not geographically proximate groups.

Taken in concert, we have argued that these types of social structure-- actor effects, reciprocity, homophily and stochastic equivalence-- influence the likelihood of violence between actors in the conflict network. Importantly, these first, second, and third order effects intuitively reflect theoretical perspectives found in the existing scholarship on civil war and rebel group behavior.  Our study unites these two literatures and demonstrates the importance of combining substantive theories of violent behavior with network estimation strategies. %This brings us to our first claim: that a model accounting for network interdependencies will better predict who fights whom and when then a purely dyadic model.

\subsection{Actor Entry}
An important and under discussed aspect of civil conflict is that actors are not static. Armed groups can enter or exit the civil conflict network at any given time. As \citet[p. 168]{kathman:wood:2015} have argued ``conflict systems are fluid, and competition varies in response to the arrival or exit of violent combat groups...", thus the entry of particularly strong or aggressive groups can have violent ripple effects across the entire conflict network. We argue that the presence of a group (call them group A) can increase the likelihood of conflict between two separate groups (group B and C). The two main ways that an actor's entry influences the behavior of other groups is by effecting the capabilities and actions of the government and altering the system-wide distribution of resources. 

The entry of a new group that clashes with the government can significantly increase the level of conflict in the network. Obviously, this is in part because the new violent group is directly involved in conflict, but this group can also spur violence between other armed groups. The entry of a new violent group that targets the government has two possible effects on the conflict network. First, the new group's entry could strain, and ultimately decrease,  government resources and capabilities. This harms the government's ability to provide security and deter violence from other armed groups. Secondly, when a new armed group successfully defeats government actors in a given region or battle, this harms the government's reputation and makes them appear more vulnerable. Thus, if the government does not want to incite other challengers \citep{walter:2006} it will try to overcompensate for it's loss and demonstrate strength by fighting heavily against armed actors. %CD: i went back and forth with saying "all armed actors" here. It sounds too big a jump, but at the same time I am not sure we clearly communicate the idea that to compensate for the threat generated by the new group, the gov will target as many other challenger groups as possible. It's actually not that intuitive, since you could also say to deal with the new scary group the gov would throw all its resources at that...

The entry of a new group can also effect the actions of other rebels based on those rebels' sources of support. In particular, the entrance of a new group might displace existing groups. This potentially leads groups into conflict with each other over scarce resources and lootable goods. At the same time, the entry of a new group might ignite a process of outbidding for the support of the civilian population who could be drawn to support the aggressive and ``successful" new group. 

The entrance of powerful or violent rebel groups into the conflict network can radically change the status quo. These groups influence perceptions and capabilities of the government and they also force other rebel groups to be increasingly aggressive in chasing scarce resources. Thus, even when we account for new groups' first order tendencies towards violence, we expect to see their entry associated with a higher level of violence in the overall system.

\emph{H1: The entrance of a new group makes violence more likely for all actors.}






% older verison: 
% In the field of interstate conflict studies, network approaches are gaining more attention.\footnote{For example, see \citet{kinne:2012,metternich:etal:2015,minhas:etal:2016}. Also see the 2016 special issues on network based approaches in International Studies Quarterly and the Journal of Peace Research.} Yet, these approaches have seen little usage in the study of civil wars even though the dynamics underlying many modern day intrastate conflicts involve interdependent interactions between multiple groups. The reason for studying dyadic interactions using a network-based approach is to acknowledge the possibility that the actions of actors in a system are contingent on one another. A significant amount of work has shown that by moving away from the assumption that interactions within a system are independent,  we can better account for the underlying process that generates relational data as well as make more accurate inferences.\footnote{For example, see \citet{%snijders:1996,
% dorff:ward:2013}.}

% Below we show a stylized version of our dependent variable of interest: the occurrence of conflict between actors in a network. Analyzing conflictual relations in a network context enables us to use information about dependencies between actors, rather than assuming these dependencies do not exist. A stylized representation of a dynamic network is shown below. In Figure~\ref{fig:netPanel2}, the nodes represent armed actors and the ties represent battles between armed actors. The graphic highlights several important features relevant to our study. Each panel represents a different phase, or new year, of the conflict. In each year, new armed actors might enter or exit the conflict (i.e., an armed group could stop fighting or dissolve). In the visualization, new entrants are shaded in grey. Second, new ties can form between actors in any period, but ties can also disappear over time. For example, moving from panel 1 to panel 2 we can see that some of the linkages in the first time period no longer appear in the second time period, which indicates that two actors that fought in the first time period are not fighting in the second time period. These features demonstrate the type of information that a network conceptualization of conflict development can provide.

% %\vspace{.5in}

% % \begin{center}
% % \textsc{Table~\ref{fig:netPanel} here}
% % \end{center}

% % \newpage
% % \begin{table}[ht]
% \begin{figure}[ht]
% 	\centering
% 	\begin{tabular}{c|c|c}
% 		$t=1$ & $t=2$ & $t=3$ \\
% 		\includegraphics[width=.2\textwidth]{netPanel2_1.pdf} &
% 		\includegraphics[width=.2\textwidth]{netPanel2_2.pdf} &
% 		\includegraphics[width=.2\textwidth]{netPanel2_3.pdf} \\
% 	\end{tabular}
% 	\caption{Stylized network graphic depicting changing connections between actors overtime, directed linkages, and community clusters.}
% 	\label{fig:netPanel2}
% % \end{table}
% \end{figure}

% % % \newpage
% % \begin{table}[ht]
% % 	\begin{tabular}{c|c|c}
% % 		$t=1$ & $t=2$ & $t=3$ \\
% % 		\includegraphics[width=.3\textwidth]{netPanel_1.pdf} &
% % 		\includegraphics[width=.3\textwidth]{netPanel_2.pdf} &
% % 		\includegraphics[width=.3\textwidth]{netPanel_3.pdf} \\
% % 	\end{tabular}
% % 	\caption{Stylized network graphic depicting changing connections between actors overtime.}
% % 	\label{fig:netPanel}
% % \end{table}

% % \newpage

% % The second iteration of 
% This visual example contains additional layers of information useful for understanding the development of conflict: directed linkages, reciprocity, and community clusters as well. By utilizing information about the direction of relationships, i.e., capturing which actors initiate battles with other actors, we can then reveal reciprocal relationships that evolve over time. The linked arrows in Figure~\ref{fig:netPanel2} demonstrates these relational dependencies. Further, we also observe a cluster of actors, or a community of armed groups, in the bottom right corner of the network -- these nodes are designated by squares instead of circles. This cluster's density increases over time and could suggest the existence of a conflict community wherein a shared latent attribute of each actor drives their propensity towards violence with one another.

% Drawing from the literature on rebel groups and non-state actors who operate in civil conflict environments, we focus on three forms of dependencies that drive the occurrence of violence between actors: actor-level effects such as actor centrality, second-order effects such as reciprocity, and third-order effects that lead certain actors to form clusters of conflict.

% \subsection*{Actor Centrality}

% We often observe heterogeneous actor centrality in social systems, revealing that some actors are more likely to initiate or receive ties (in this case, battles). Heterogeneity across actor centrality provides us with information on which set of events led to the creation of the observed network \citep{barabasi:reka:1999}.  Heterogeneity \textit{across} actors' activity in a network produces homogeneous ties \textit{between} actors. For example, \cite{weinstein:2007} argues that rebel groups with more ideologically motivated followers are thought to be less violent than groups whose members are motivated by greed. The motivation of group members, then, is a latent, unobserved attribute that accounts for why one rebel group is more prone to attack other groups. Thus, the relationships between the set of dyads $\{i \rightarrow j,\, i \rightarrow k,\, i \rightarrow l\}$ will be more similar to each other than they are to other interactions in the system, such as $m \rightarrow n$. This occurs because each of those interactions involve actor $i$ as the sender of the event \citep{kenny:lavoie:1984}, and sender $i$ possesses some latent attribute that makes the group more conflictual.

% %Another straightforward example of why this type of dependence emerges in intrastate conflict networks involves government actors. Government actors likely have interests in asserting their control across the country. Specifically, in a scenario where rebel groups $\{j,k,l\}$ operate in different parts of a country, we would expect government actor $i$ to be more likely to initiate conflict with those rebel groups than an actor with non-competing interests. This logic naturally extends to non-government groups. For example, a rebel group may act upon ambitions of establishing control over broad swaths of the country, and, as a result, become a central initiator of conflicts in the network. For similar reasons we may observe heterogeneity in how likely actors are to receive conflict events, and, relatedly, find that actors initiating more battles are also likely to appear on the receiving end of violent attacks. Each of these processes creates dependence across observations that involve the same actor. 

% \subsection*{Reciprocity}

% Reciprocity is the notion that the interaction between $i \rightarrow j$ and $j \rightarrow i$ are dependent upon one another. The concept of reciprocity has deep roots in the study of relations between states \citep{richardson:1960}%,keohane:1989},
% and it is no less relevant for studying relations between actors in an intrastate conflict context. Reciprocity captures the intuition that if one actor is, for example, particularly aggressive towards another actor in the network, the recipient of such aggression will ``respond in kind'' and exhibit aggressive behavior in return. In addition to this direct relationship, where attacks lead to retaliation, we might also see reciprocity when a common factor is both causing $i$ to attack $j$ and $j$ to attack $i$. For example, literature on competition between armed groups suggests that when two groups are competing over scarce resources, \citep{fjelde:nilsson:2012}, we should expect each group to attack the other and thus attacks by $i$ on $j$ are likely to be paired with attacks by $j$ on $i$. We expect that in the context of battles, reciprocity is likely to play a role in motivating actors to respond to attacks by opponent groups.

% \subsection*{Homophily \& Stochastic Equivalence}

% Apart from these nodal and dyadic patterns that commonly arise in social systems, networks often exhibit patterns of dependence involving multiple actors. These types of patterns can broadly be summarized in terms of the concepts of stochastic equivalence and homophily \citep{wasserman:faust:1994}. Stochastic equivalence may manifest in a network in which actors cluster into groups, and the group determines an actor's patterns of relations with other groups. More concretely, a pair of actors $ij$ are stochastically equivalent if the probability of $i$ relating to, and being related to, by every other actor is the same as the probability for $j$ \citep{anderson:etal:1992}. Consider a case in which actors nest in groups based on, for example, a shared but unobserved ideology: $\mathbf{a}=\{i,j\}$ and $\mathbf{b}=\{k,l\}$. Here we may likely find that the relations of $i$ towards $k$ and $l$ is similar to how $j$ behaves towards those actors, and this induces a dependence in the relations of actors within the system. Specifically, the behavior of actors within the conflict system may be a function of a latent community to which they belong. A large body of literature in conflict studies motivates this reasoning, since ethnicity and religion are often tied to the behavior of actors within their community context \citep{reynalquerol:2002, denny:walter:2014}.  In the case of the conflict in Nigeria, we might see groups that share a common religion attacking a group whose followers follow a different religion, or groups with a given tribal identity coming into conflict with a group from a rival tribe.

% Homophily is a type of dependence pattern involving multiple actors that often leads to the emergence of clusters within social systems \citep{shalizi:thomas:2011}. In the intrastate setting, conflict clusters may emerge when multiple actors are competing for control over the same territory. Typically, conflict scholars attempt to account for this by explicitly identifying a territorial objective, but in some cases this may not be observed. Using a network based approach, we can examine the conflictual patterns of actors to parse out this underlying geographic structure. Thus a latent shared attribute that explains the occurrence of homophily in relational data may be geographical in nature.
% % \footnote{A clique in a  network is a structure in which three actors $\{i,j,k\}$ each interact with one another to form a triangle.}
% % \footnote{Interestingly, \citet{wong:etal:2006} show that homophily in social networks is frequently determined by geographical space.} 

% The principal idea underlying each of the dependence patterns discussed above is that the actions, or inactions, of any actor, are taking place within an interdependent system. The interdependencies result because of latent attributes possessed by a single rebel group or shared by many groups. Attempting to estimate the presence of these interdependencies allows us to better understand whether there is latent structure to an intrastate conflict system. Further by incorporating the estimation of these network effects into a modeling framework, we reduce the likelihood of making faulty inferences and improve our ability to forecast conflict.

% \subsection*{Entrance of Key Players}

% An often undiscussed concept in both networks and conflicts relates to the ripple effects that the entrance of a key player can have on a system. The entry of a particularly active or violent group could lead to an increase in conflict through a number of mechanisms. As \citet[p. 168]{kathman:wood:2015} have argued ``conflict systems are fluid, and competition varies in response to the arrival or exit of violent combat groups...", thus the entrance of certain groups could increase levels of conflict. One possible mechanism through which violence might increase is through a process of rebel group outbidding, where in order for different groups to get support and resources from the population, they need to prove their conviction and capabilities through increasingly violent acts \citep{bloom:2004, nemeth:2014}. The entrance of additional groups might also make ending conflict harder, as \citet{cunningham:etal:2009} notes, because the entrance of an additional (dominant) group makes de-escalating conflict more difficult as it increases the number of veto players in peace negotiations.

% We suggest that the entrance of key players affects violence because of their strategic relationship with the government. The arrival of a new challenger group that successfully challenges the government may shift perceptions of the government as weak or vulnerable, leading other groups to more willingly challenge the government, as shown in \citet{walter:2006} or challenge other armed groups, as studied by \citet{fjelde:nilsson:2012}. This particular dynamic also incentivizes the government to react with violence, rather than negotiation in order to deter the development of future challengers.  \citet[p. 613]{fjelde:nilsson:2012} find that ``states with weak coercive power create opportunities for nonstate actors to engage in armed struggle against each other.''

% Boko Haram is an excellent example of this effect in Nigeria. Boko Haram gained such success that in mid-2014 they were able to effectively control swathes of territory in and around their home state of Borneo. The rise of Boko Haram and the challenges it posed to Nigeria's security apparatus provides incentives for other groups to increase their violent behavior, even in dyads that do not involve Boko Haram.
